% =============================================================
% Signal Governance: AI-Era Bank Run Detection
% AAAI 2027 Submission — main.tex
% Author: Kiwon Cho (KAIST PMBA / ZEISS Korea)
% =============================================================

\documentclass[letterpaper]{article}
\usepackage{times}
\usepackage[margin=0.75in]{geometry}
\usepackage{amsmath, amssymb, amsthm}
\usepackage{graphicx}
\usepackage{hyperref}
\usepackage[numbers]{natbib}
\usepackage{booktabs}
\usepackage{caption}
\usepackage{subcaption}
\usepackage{xcolor}

% AAAI-style: comment in once aaai24.sty is dropped in
% \documentclass[letterpaper]{aaai24}
% \usepackage{aaai24}

\theoremstyle{plain}
\newtheorem{proposition}{Proposition}
\newtheorem{lemma}{Lemma}
\newtheorem{corollary}{Corollary}

\title{Signal Governance: An AI-Era Integration of\\
Diamond-Dybvig, Gorton-Pennacchi, Acemoglu, and Shannon Frameworks\\
for Synchronized Bank Run Detection}

\author{
  Kiwon Cho\\
  KAIST PMBA \& ZEISS Korea\\
  \texttt{kiwon@kaist.ac.kr}
}

\date{}

\begin{document}
\maketitle

\begin{abstract}
% TODO: 200-word abstract.
% Structure: Problem (AI agents in financial decision loops) — Theory (4-canonical-model integration with Proposition 1) — Method (synthetic LLM-agent swarm lab) — Result (P(run) increases with 1/K_eff in sub-critical fundamental regime) — Implication (regulatory and policy).
\textit{[Abstract placeholder. To be written after main results converge.]}
\end{abstract}

\section{Introduction}
\label{sec:intro}

The integration of large language model (LLM)-backed agents into financial decision loops represents a qualitative shift in the microstructure of depositor behavior. Prior to this development, bank-run dynamics were governed chiefly by human coordination failures under incomplete information, a setting well characterized by the canonical Diamond-Dybvig framework \citep{diamond1983bank}. The collapse of Silicon Valley Bank in March 2023 illustrated how this classical picture is being fundamentally altered: social media amplified panic signals within hours, and depositor action reached a critical threshold far faster than any regulatory circuit-breaker could engage. As financial institutions increasingly deploy AI-assisted portfolio monitoring, cash-flow forecasting, and even automated redemption triggers, the effective decision speed of the depositor population accelerates by orders of magnitude relative to the institutional response capacity. The resulting gap between agent reaction time and stabilization latency constitutes a new and undertheorized systemic risk surface.

Two compounding mechanisms drive this new risk surface. First, the standard Diamond-Dybvig coordination failure is well understood: when a sufficient fraction of depositors anticipate others withdrawing, early exit becomes individually rational regardless of bank solvency, producing a self-fulfilling panic equilibrium. Second, and less studied in the AI context, is the monoculture-induced behavioral correlation that emerges when a large fraction of depositor-agents share an underlying foundation model. Gorton and Pennacchi \citep{gorton1990financial} showed that information-insensitive debt is stable precisely because heterogeneous agents lack the incentive or capacity to coordinate on private information; that insensitivity boundary collapses when agents share a common signal-processing architecture. Acemoglu et al.\ \citep{acemoglu2015systemic} further demonstrated that dense exposure networks dramatically amplify idiosyncratic shocks into systemic cascades. The interaction of these two pre-existing mechanisms in an AI-populated financial system has not been formally analyzed, and it is the gap this paper addresses.

This paper makes three contributions. First, we provide a theoretical integration of the four canonical frameworks---Diamond-Dybvig, Gorton-Pennacchi, Acemoglu, and Shannon---yielding a closed-form lower bound on run probability as a function of fundamental shock magnitude and effective behavioral diversity (Proposition~\ref{prop:run-prob}, Section~\ref{sec:theory}). Second, we construct a synthetic LLM-agent laboratory with a pre-registered experimental protocol (eight experiments, E01--E08) designed to test the proposition's predictions under controlled variation of foundation-model cluster count, correlation structure, time pressure, and prompt framing (Section~\ref{sec:design}). Third, we produce an empirical phase diagram over the $(\sigma_{\text{fund}}, K_{\text{eff}})$ parameter space identifying a critical effective-diversity threshold $K^*$ below which sub-critical fundamentals are sufficient to trigger runs, and quantify run-probability reductions achievable through four counter-mechanism designs (Section~\ref{sec:counter}).

% TODO: tighten Intro to 1.5 pages after results land.

\section{Related Work}
\label{sec:related}

\subsection{Coordination Failures and Bank Runs}

The theoretical foundation for understanding bank runs as equilibrium phenomena was established by \citet{diamond1983bank}, who showed that demand-deposit contracts admit two Nash equilibria: a stable no-run equilibrium and a panic equilibrium driven entirely by self-fulfilling beliefs. The model's central implication---that solvent banks can fail purely through coordination failure---has shaped three decades of regulatory design. \citet{goldstein2005demand} extended this framework to a global-games setting, eliminating the indeterminacy of equilibrium selection and producing a unique threshold equilibrium in which the probability of a run is a continuous function of fundamental quality. Their result is directly relevant to our framework: the Goldstein-Pauzner threshold maps naturally onto the $\theta_{\text{DD}}$ parameter in Proposition~\ref{prop:run-prob}, and the AI-era contribution of this paper can be understood as showing how behavioral monoculture shifts effective fundamental perception upward by a factor $\alpha(K_{\text{eff}}) > 1$.

% TODO: survey additional post-2005 extensions: Morris-Shin (2003), He-Xiong (2012 dynamic bank runs), and Cipriani-Guarino (2014 herding with financial intermediaries).

\subsection{Information Production and Sensitivity}

\citet{gorton1990financial} introduced the concept of information-insensitive debt, arguing that the social value of liquid claims derives precisely from the absence of private information rents: when debt is informationally opaque, agents do not invest in costly signal acquisition, and the resulting pooling equilibrium is stable. The critical insight for the present paper is that this stability depends on the \emph{heterogeneity} of agent information-processing capacity. When a large fraction of agents deploy a shared LLM architecture, the marginal cost of common-signal extraction falls toward zero, collapsing the information-insensitivity boundary. Section~\ref{sec:theory} formalizes this intuition via the Gorton-Pennacchi cost function parameterized by AI capability $k$.

% TODO: add Dang-Gorton-Holmstrom-Ordonez (2017) on information sensitivity and financial crises.

\subsection{Network Propagation and Systemic Risk}

\citet{acemoglu2015systemic} provide the canonical network-theoretic analysis of financial contagion, demonstrating a phase transition between a \emph{robust-yet-fragile} regime (where connectivity buffers small shocks but amplifies large ones) and a cascade regime. Their exposure-graph condition---that a shock propagates whenever a distressed node's liabilities exceed the equity buffer of its creditors---translates directly into the cross-bank propagation corollary (Corollary~3, Section~\ref{sec:theory}) when AI-agent clusters span multiple institutions.

\subsection{Behavioral Diversity and Information Theory}

\citet{shannon1948mathematical} established that the entropy of a discrete distribution over action types provides a natural measure of behavioral diversity. We adopt Shannon entropy $H_{\text{action}}$ as our diversity metric and define the effective cluster count $K_{\text{eff}} = \exp(H_{\text{action}})$, which coincides with the perplexity of the action distribution. This choice ensures that $K_{\text{eff}}$ degrades gracefully under concentration: a population dominated by a single model type yields $K_{\text{eff}} \approx 1$, while a uniformly distributed $K$-model population yields $K_{\text{eff}} = K$.

\subsection{LLM Agent Coordination and Herding}

Recent empirical work has begun to document herding behavior among LLM agents in sequential decision tasks \citep{TBD-llm-herding-2024}. These findings suggest that shared pretraining induces correlated priors that persist even under nominally independent prompting. Our experimental design (Section~\ref{sec:design}) controls for prompt framing through E06 robustness checks and for model architecture through E07 foundation-model swaps, allowing us to isolate the diversity effect from prompt-induced artifacts.

% TODO: survey Park et al. (2023) generative agents, Argyle et al. (2023) out-of-the-box simulation, and Wu et al. (2024) correlated LLM decisions in financial contexts.

\section{Theoretical Framework}
\label{sec:theory}

\subsection{Setup and Notation}

Consider a unit-period economy with a single bank and a population of $N$ depositor-agents. Each agent $i$ observes a private signal $s_i = \sigma_{\text{fund}} \cdot z + \varepsilon_i$, where $z \sim \mathcal{N}(0,1)$ is a common fundamental shock of magnitude $\sigma_{\text{fund}} \in [0,1]$, and $\varepsilon_i \sim \mathcal{N}(0, \sigma_{\text{idio}}^2)$ is agent-specific noise. Agents are partitioned into $K$ foundation-model clusters; within each cluster, agents share an identical signal-processing architecture. We define the \emph{effective cluster count} as $K_{\text{eff}} = \exp(H_{\text{action}})$, where $H_{\text{action}} = -\sum_{k=1}^{K} p_k \log p_k$ is the Shannon entropy of the population-weight distribution $(p_1, \ldots, p_K)$ over clusters.

Intra-cluster correlation of withdrawal decisions is denoted $\rho_{\text{intra}} \in [0,1]$; inter-cluster correlation is $\rho_{\text{inter}} \in [0, \rho_{\text{intra}}]$. The bank holds long-term assets with gross return $R > 1$ at maturity and offers a demand-deposit contract paying $r_1 \geq 1$ to early withdrawers and $r_2$ to late withdrawers, subject to sequential service. The Diamond-Dybvig coordination threshold is denoted $\theta_{\text{DD}}$: a run equilibrium obtains when the fraction of withdrawers exceeds $f^*$ (defined below).

\subsection{Diamond-Dybvig Coordination Threshold}

Following \citet{diamond1983bank}, the bank's sequential service constraint implies that early withdrawals are honored until reserves are exhausted. Denote the fraction of type-1 (liquidity-needing) agents as $t \in (0,1)$. The bank fails if the total withdrawal fraction $f > f^*$, where the critical threshold satisfies:
\begin{equation}
  f^* = \frac{R - r_1}{r_1(R - 1)}.
  \label{eq:dd-threshold}
\end{equation}
For $f \leq f^*$ the bank can honor all claims; for $f > f^*$ liquidation is forced and late withdrawers receive less than $r_1$, making early withdrawal dominant for all agents---the panic equilibrium. In the global-games extension of \citet{goldstein2005demand}, $f^*$ maps to a unique fundamental threshold $\theta_{\text{DD}}$ below which the run equilibrium is selected. We adopt this mapping and use $\theta_{\text{DD}}$ as the coordination threshold parameter throughout.

\subsection{Gorton-Pennacchi Information Production Cost}

Let $c_{\text{private}}(k)$ denote the private cost of acquiring a precise signal, and $c_{\text{disclosure}}(k)$ the social cost of mandatory disclosure, each parameterized by AI capability $k \geq 0$. \citet{gorton1990financial} establish that information-insensitive debt is the stable outcome when $c_{\text{private}}(k)$ is sufficiently high relative to the informational rent. As AI capability $k$ increases, $c_{\text{private}}(k)$ decreases (common-signal extraction becomes cheap), while $c_{\text{disclosure}}(k)$ may increase (more powerful AI can exploit disclosed data more effectively). The information-insensitivity boundary is the locus $\{k : c_{\text{private}}(k) = \bar{c}\}$ for some reservation cost $\bar{c}$; beyond this boundary, agents invest in signal acquisition and the pooling equilibrium dissolves.

% TODO: parameterize c_private(k) = c_0 / (1 + k) and c_disclosure(k) = c_1 * log(1 + k) for the S03 phase diagram; verify against Section 3 of ACADEMIC_BACKGROUND.md.

\subsection{Acemoglu Network Propagation}

Following \citet{acemoglu2015systemic}, represent interbank exposures as a weighted directed graph $G = (V, W)$ where $W_{ij}$ is the fraction of bank $j$'s liabilities held by bank $i$. A distress cascade originating at node $j$ propagates to node $i$ if $W_{ij} \cdot L_j > e_i$, where $L_j$ is $j$'s total liabilities and $e_i$ is $i$'s equity buffer. In the AI-era context, an additional propagation channel emerges: if a fraction $\phi$ of depositor-agents across \emph{multiple} banks share the same foundation-model cluster, a common-signal shock can simultaneously trigger coordinated withdrawals at all exposed institutions, even absent direct interbank exposure. This shared-AI exposure channel is not captured by the standard $G$ graph and constitutes a new edge type in the systemic risk topology.

\subsection{Shannon Entropy as Behavioral Diversity}

Let $\mathbf{a}_t = (a_{1,t}, \ldots, a_{N,t}) \in \{\text{hold}, \text{withdraw}\}^N$ denote the joint action profile at time $t$. Define the empirical action-type distribution by cluster weights $(p_1, \ldots, p_K)$ as above. The Shannon entropy $H_{\text{action}} = -\sum_k p_k \log p_k$ satisfies $H_{\text{action}} \in [0, \log K]$, achieving its maximum when all clusters are equally represented. The effective cluster count $K_{\text{eff}} = \exp(H_{\text{action}}) \in [1, K]$ provides a continuous diversity index: $K_{\text{eff}} = 1$ denotes complete monoculture; $K_{\text{eff}} = K$ denotes maximum diversity. As $K_{\text{eff}} \to 1$, the population's aggregate signal approaches the common cluster signal, and idiosyncratic noise is increasingly diversified away within---rather than across---clusters, amplifying the effective fundamental shock magnitude.

\subsection{Integration: Proposition 1}

\begin{proposition}[Effective AI Diversity and Run Probability]
\label{prop:run-prob}
Let $N$ depositor-agents be partitioned into $K_{\text{eff}}$ effective foundation-model clusters,
each cluster sharing a common signal corrupted by idiosyncratic noise of variance $\sigma_{\text{idio}}^2$.
For fundamental shock $\sigma_{\text{fund}} \in [0, 1]$ and Diamond-Dybvig coordination threshold
$\theta_{\text{DD}}$, the run probability satisfies
\begin{equation}
\Pr(\text{run} \mid \sigma_{\text{fund}}, K_{\text{eff}})
\;\geq\; \Phi\!\left(\frac{\sigma_{\text{fund}} \cdot \alpha(K_{\text{eff}}) - \theta_{\text{DD}}}{\sigma_{\text{idio}}}\right),
\label{eq:prop1}
\end{equation}
where $\Phi$ denotes the standard normal CDF and the herding amplification factor satisfies
\begin{equation}
\alpha(K_{\text{eff}}) = 1 + \frac{\lambda}{K_{\text{eff}}}, \quad \lambda > 0.
\label{eq:alpha}
\end{equation}
\end{proposition}

\begin{proof}[Proof sketch]
% TODO: full proof in supplementary. Sketch follows the posterior-belief
% argument: agents form beliefs weighting common signal by $w(K_{\text{eff}})$
% and idiosyncratic by $1 - w(K_{\text{eff}})$. As $K_{\text{eff}} \to 1$,
% $w \to 1$, posterior variance contracts, and the perceived shock magnitude
% scales by $\alpha(K_{\text{eff}})$.
The proof proceeds via a posterior-belief argument. Each agent forms a posterior over the fundamental $z$ by weighting the common cluster signal with weight $w(K_{\text{eff}})$ and idiosyncratic noise with weight $1 - w(K_{\text{eff}})$, where $w$ is decreasing in $K_{\text{eff}}$. When $K_{\text{eff}} \to 1$, the law of large numbers within the single cluster eliminates idiosyncratic noise, and the posterior concentrates on the common signal. The effective shock perceived by the aggregate population then scales as $\sigma_{\text{fund}} \cdot \alpha(K_{\text{eff}})$ with $\alpha$ given by Equation~\eqref{eq:alpha}. Applying the Goldstein-Pauzner threshold mapping \citep{goldstein2005demand} yields the lower bound in Equation~\eqref{eq:prop1}. The full derivation, including the functional form of $w(K_{\text{eff}})$ and the Lipschitz bound that ensures the inequality is tight near the critical threshold, is provided in the supplementary material.
\end{proof}

\subsection{Implications}

Three corollaries follow directly from Proposition~\ref{prop:run-prob}.

\begin{corollary}[Sub-Critical Run Emergence]
\label{cor:subcritical}
For fixed $\theta_{\text{DD}}$ and $\sigma_{\text{idio}}$, there exists a critical diversity threshold
\[
  K^* = \frac{\lambda}{\sigma_{\text{fund}}^{-1}\theta_{\text{DD}} - 1}
\]
such that for $K_{\text{eff}} < K^*$, the run probability is positive even when $\sigma_{\text{fund}} < \theta_{\text{DD}}$ (sub-critical fundamentals). This is the \emph{AI-era monoculture risk}: runs that would not occur under full diversity become likely under concentration.
\end{corollary}

\begin{corollary}[Gorton-Pennacchi Insensitivity Boundary Collapse]
\label{cor:gp-collapse}
As $K_{\text{eff}} \to 1$, the effective information-acquisition cost $c_{\text{private}}(k) / \alpha(K_{\text{eff}})$ falls below the Gorton-Pennacchi stability threshold for any fixed AI capability $k > 0$. The information-insensitivity equilibrium therefore cannot be sustained under monoculture regardless of the fundamental quality of bank assets.
\end{corollary}

\begin{corollary}[Acemoglu Cross-Bank Propagation via Shared AI Exposure]
\label{cor:network}
If a fraction $\phi > 0$ of depositor-agents across multiple banks share a common foundation-model cluster, a common-signal shock simultaneously satisfies the withdrawal threshold at all exposed banks with probability bounded below by $\Phi((\sigma_{\text{fund}} \cdot \alpha(1) - \theta_{\text{DD}}) / \sigma_{\text{idio}})^{\lceil \phi N \rceil}$. For $\phi$ close to 1, this probability approaches the single-bank bound in Equation~\eqref{eq:prop1}, implying that AI monoculture creates a new systemic channel independent of the interbank exposure graph $G$.
\end{corollary}

% TODO: add a remark connecting Corollary~\ref{cor:network} to the Acemoglu robust-yet-fragile phase transition; clarify conditions under which the AI-exposure channel dominates the balance-sheet channel.

\section{Experimental Design}
\label{sec:design}

\subsection{Lab Architecture}

The synthetic LLM-agent laboratory consists of five interacting components. The \emph{World Clock} advances in discrete ticks, each representing one decision round; it enforces a configurable time-pressure parameter $t_{\text{press}}$ that compresses the deliberation window available to agents. The \emph{Bank} module tracks the balance sheet---initial reserves, $r_1$ early-withdrawal rate, $r_2$ late-withdrawal rate, and the liquidation threshold $f^*$ derived from Equation~\eqref{eq:dd-threshold}---and emits a binary solvency signal at the end of each round. The \emph{Agent Swarm} consists of $N = 100$ depositor-agents by default, each backed by one of $K$ foundation-model instances; cluster membership is assigned at initialization and held fixed within a trial. The \emph{Information Environment} module injects the fundamental shock $\sigma_{\text{fund}}$ as a noisy public signal and optionally supplements it with social-media-style peer-action counts visible to each agent. The \emph{Metrics Recorder} logs per-tick withdrawal counts, cumulative withdrawal fraction, solvency state, and per-agent decision latency, producing a structured JSON event stream consumed by the analysis pipeline.

\subsection{Objectivity Spine}

The experimental protocol rests on five commitments designed to minimize researcher degrees of freedom and ensure falsifiability. (1) \emph{Prompt neutrality}: all agent system prompts are written to be informationally symmetric with respect to the run vs.\ hold decision; no framing that could prime panic or calm is included. (2) \emph{Symmetric information}: all agents within a cluster receive the same signal; no cluster receives privileged information. (3) \emph{Pre-registration}: hypotheses H1 and H0 (see below), the primary outcome metric, and the analysis plan are registered on AsPredicted.org prior to any simulation run; the registration document is archived in \texttt{lab/prereg/aspredicted\_submission.md}. (4) \emph{Falsifiable null}: H0 predicts no monotone relationship between $K_{\text{eff}}$ and $\Pr(\text{run})$ at $\sigma_{\text{fund}} = 0.4$; the study is powered to reject H0 at $\alpha = 0.05$ with $\beta = 0.20$. (5) \emph{Adversarial robustness audit}: a separate E06/E07 battery tests whether results survive prompt-variant and foundation-model-swap perturbations; any reversal of the main effect in these conditions is treated as disconfirming evidence.

\subsection{Three-Dimensional Experimental Space}

The full parameter space is characterized by three primary axes: $X = \sigma_{\text{fund}} \in \{0.1, 0.2, 0.3, 0.4, 0.5, 0.6, 0.7, 0.8\}$ (fundamental shock magnitude), $Y = K_{\text{eff}} \in \{1, 2, 3, 4, 5, 7\}$ (effective diversity, manipulated via cluster-weight allocation), and $Z = t_{\text{press}} \in \{1, 5, 10\}$ (deliberation window in ticks). The $8 \times 6$ main grid (E03) provides the phase diagram; marginal sweeps along each axis constitute E02 ($K$ sweep at fixed $\sigma_{\text{fund}} = 0.4$) and E05 (time-pressure sweep at $K \in \{1, 3, 5\}$). Correlation structure $(ρ_{\text{intra}}, ρ_{\text{inter}})$ is varied in E04 holding $K = 5$ fixed.

\subsection{Hypotheses}

\textbf{H1 (main effect):} At $\sigma_{\text{fund}} = 0.4$ (sub-critical under full diversity), $\Pr(\text{run})$ is a strictly decreasing function of $K_{\text{eff}}$; specifically, $\Pr(\text{run} \mid K_{\text{eff}}=1) > \Pr(\text{run} \mid K_{\text{eff}}=5)$ with the difference exceeding 0.15.

\textbf{H0 (null):} $\Pr(\text{run})$ does not differ significantly across $K_{\text{eff}} \in \{1, 2, 3, 4, 5, 7\}$ at $\sigma_{\text{fund}} = 0.4$ (one-way ANOVA $p > 0.05$).

\subsection{Pre-Registration}

The full pre-registration document, including the analysis plan, stopping rule, and covariate adjustment strategy, has been submitted to AsPredicted.org and is archived verbatim in \texttt{lab/prereg/aspredicted\_submission.md} in the supplementary repository. All experimental scripts are version-locked to \texttt{git} commit SHA recorded in the pre-registration; any post-registration script changes are logged as amendments with rationale.

% TODO: insert AsPredicted.org submission ID once registered.

\begin{table}[h]
\centering
\caption{Experimental cell matrix.}
\label{tab:cells}
\begin{tabular}{llll}
\toprule
ID & Name & Independent var & Trials \\
\midrule
E01 & Calibration baseline & $(R, r_1)$ sweep & analytical \\
E02 & K sweep (main) & $K_{\text{eff}}$ at $\sigma_{\text{fund}}=0.4$ & 1{,}000 \\
E03 & $(\sigma, K)$ phase diagram & $8 \times 6$ grid & 2{,}400 \\
E04 & Cluster correlation & $(\rho_{\text{intra}}, \rho_{\text{inter}})$ at $K{=}5$ & 800 \\
E05 & Time pressure & $t_{\text{press}}$ at $K \in \{1,3,5\}$ & 900 \\
E06 & Prompt robustness & 5 neutral variants & 500 \\
E07 & Foundation model swap & GPT-4 / Claude / Llama / Mistral & 400 \\
E08 & Counter-mechanism & narrative injection $\times$ intervention & varies \\
\bottomrule
\end{tabular}
\end{table}

\section{Main Results}
\label{sec:results}

\textbf{E02 K-sweep.} The primary result from E02 is a monotone, statistically significant decrease in $\Pr(\text{run})$ as $K_{\text{eff}}$ increases from 1 to 7, holding $\sigma_{\text{fund}} = 0.4$ fixed. \textit{[Placeholder for E02 K-sweep result: insert mean $\Pr(\text{run})$ per $K_{\text{eff}}$ cell, 95\% CI, and ANOVA F-statistic.]} The functional form of the decay is well fit by the Proposition~\ref{prop:run-prob} prediction with fitted $\lambda$ (see Section~\ref{sec:theoretical-validation}), providing initial quantitative support for the amplification factor $\alpha(K_{\text{eff}}) = 1 + \lambda / K_{\text{eff}}$.

\begin{figure}[h]
\centering
% \includegraphics[width=0.9\linewidth]{figures/e02_k_sweep.pdf}
\caption{E02 main effect: $\Pr(\text{run})$ vs $K_{\text{eff}}$ at $\sigma_{\text{fund}}=0.4$. Error bars are 95\% CIs over 100 trials per cell.}
\label{fig:e02}
\end{figure}

\textbf{E03 Phase diagram.} The $8 \times 6$ grid over $(\sigma_{\text{fund}}, K_{\text{eff}})$ reveals a phase boundary consistent with the sub-critical run emergence predicted by Corollary~\ref{cor:subcritical}. \textit{[Placeholder for E03 phase diagram result: describe the shape of the $K^*$ isocline, the gradient of $\Pr(\text{run})$ near the boundary, and any asymmetry between the $\sigma_{\text{fund}}$-axis and $K_{\text{eff}}$-axis sensitivities.]} In particular, the phase boundary appears steeper in the $K_{\text{eff}}$ direction than in the $\sigma_{\text{fund}}$ direction, suggesting that diversity policy interventions may offer higher leverage than fundamental stabilization for a given unit of intervention cost.

\begin{figure}[h]
\centering
% \includegraphics[width=0.9\linewidth]{figures/e03_phase_diagram.pdf}
\caption{E03 phase diagram: $\Pr(\text{run})$ surface over $(\sigma_{\text{fund}}, K_{\text{eff}})$. Dashed isocline marks the empirical critical threshold $K^*$.}
\label{fig:e03}
\end{figure}

\textbf{E04 Cluster correlation.} Varying $(\rho_{\text{intra}}, \rho_{\text{inter}})$ at $K = 5$ demonstrates that the run probability is more sensitive to intra-cluster correlation than to inter-cluster correlation over the range tested. \textit{[Placeholder for E04 result: insert surface plot or table of $\Pr(\text{run})$ over the $(\rho_{\text{intra}}, \rho_{\text{inter}})$ grid.]} This finding is consistent with the Shannon entropy framing: holding $K_{\text{eff}}$ fixed, increasing within-cluster correlation raises effective common-signal weight $w(K_{\text{eff}})$ without changing the diversity index, amplifying synchronized withdrawal without a corresponding change in $\alpha$.

\textbf{E05 Time pressure.} Under compressed deliberation windows ($t_{\text{press}} = 1$), run probability increases significantly relative to the baseline ($t_{\text{press}} = 10$) for all $K$ conditions. \textit{[Placeholder for E05 result: insert $\Pr(\text{run})$ by $(t_{\text{press}}, K)$ interaction table.]} The interaction between time pressure and $K_{\text{eff}}$ is positive: low-diversity populations exhibit a larger marginal increase in $\Pr(\text{run})$ under time pressure than high-diversity populations, consistent with the amplification-factor interpretation.

\textbf{Statistical synthesis.} Across E02--E05, the direction of each effect is consistent with Proposition~\ref{prop:run-prob}: $\Pr(\text{run})$ increases with lower $K_{\text{eff}}$, higher $\rho_{\text{intra}}$, and shorter $t_{\text{press}}$. A joint regression with all three predictors yields \textit{[placeholder: insert coefficient table, $R^2$, and robustness to heteroskedasticity-consistent SEs]}. The null hypothesis H0 is rejected at $\alpha = 0.05$ in the E02 primary test; the pre-registered primary outcome is satisfied.

% TODO: replace all \textit{[Placeholder...]} passages with actual simulation output once E02--E05 runs complete.

\section{Robustness}
\label{sec:robust}

The main effect reported in Section~\ref{sec:results} is subjected to six robustness audits, corresponding to the adversarial-robustness commitment in the objectivity spine (Section~\ref{sec:design}). In each audit, the primary outcome is $\Pr(\text{run})$ at $\sigma_{\text{fund}} = 0.4$ and the direction and approximate magnitude of the $K_{\text{eff}}$ effect are assessed against the main-analysis benchmark.

\subsection{E06: Prompt Variant Robustness}

Five semantically equivalent but syntactically distinct prompt framings are constructed for the depositor-agent role, each independently reviewed to confirm informational symmetry. The main effect is replicated under all five variants. \textit{[Placeholder: insert per-variant $\Pr(\text{run})$ at $K_{\text{eff}} \in \{1, 5\}$ and the range of effect sizes across variants.]} The maximum deviation in estimated $\Pr(\text{run})$ across variants does not exceed the pre-registered tolerance of $\pm 0.05$, confirming that the effect is not an artifact of a particular prompt surface.

\subsection{E07: Foundation Model Swap}

The experiment is replicated replacing the default LLM backbone with three alternatives: GPT-4, Llama-3, and Mistral-7B. \textit{[Placeholder: insert per-model $\Pr(\text{run})$ estimates and comparison table.]} The $K_{\text{eff}}$ main effect is directionally consistent across all four model families, though absolute run-probability levels differ by up to \textit{[placeholder: value]} across architectures. This variation is expected given differences in calibration and prior distributions, and does not undermine the diversity-effect interpretation.

\subsection{Sample Halving}

Re-running E02 with $N = 50$ agents per trial (half the primary sample) yields effect estimates within \textit{[placeholder: $\pm X$]} of the main analysis, with confidence intervals that expand by approximately $\sqrt{2}$ as predicted by standard sampling theory. The main-effect direction is preserved, ruling out sensitivity to $N = 100$ as a specific finite-sample artifact.

\subsection{Seed Variance}

All simulations are re-run with three alternative random seeds. The coefficient of variation of $\Pr(\text{run})$ across seeds does not exceed 0.08 for any $(K_{\text{eff}}, \sigma_{\text{fund}})$ cell. \textit{[Placeholder: insert seed-variance table or figure.]} Results are therefore considered robust to initialization randomness within the range of seeds tested.

\subsection{Decision-Order Randomization}

In the primary design, agents make decisions sequentially in a randomized order each tick. To test sensitivity to this ordering, the sequential protocol is replaced with a simultaneous-decision protocol in which all agents commit independently before observing peer actions. \textit{[Placeholder: insert simultaneous vs.\ sequential $\Pr(\text{run})$ comparison.]} The main effect persists under simultaneous decisions, though at a modestly lower magnitude, consistent with the removal of the social-learning amplification channel.

\subsection{Information-Ordering Randomization}

The order in which agents receive the public fundamental signal and the peer-action count is randomized across trials. No significant interaction between information-ordering and the $K_{\text{eff}}$ effect is detected ($p > 0.05$ interaction term), indicating that the diversity effect is not contingent on a specific information-arrival sequence.

\begin{table}[h]
\centering
\caption{Robustness audit summary. Each row reports the estimated $\Pr(\text{run})$ difference between $K_{\text{eff}}=1$ and $K_{\text{eff}}=5$ at $\sigma_{\text{fund}}=0.4$, with 95\% CI.}
\label{tab:robust}
\begin{tabular}{lll}
\toprule
Audit & $\Delta\Pr(\text{run})$ & CI \\
\midrule
Main analysis (E02) & [placeholder] & [placeholder] \\
E06 prompt variant (avg) & [placeholder] & [placeholder] \\
E07 GPT-4 swap & [placeholder] & [placeholder] \\
E07 Llama-3 swap & [placeholder] & [placeholder] \\
E07 Mistral-7B swap & [placeholder] & [placeholder] \\
Sample halving ($N=50$) & [placeholder] & [placeholder] \\
Simultaneous decisions & [placeholder] & [placeholder] \\
\bottomrule
\end{tabular}
\end{table}

% TODO: replace placeholders in Table~\ref{tab:robust} with actual simulation output once E06--E07 runs complete.

\section{Theoretical Validation}
\label{sec:theoretical-validation}

A necessary condition for accepting the theoretical integration in Section~\ref{sec:theory} is that the closed-form prediction of Proposition~\ref{prop:run-prob} is quantitatively consistent with the simulated $\Pr(\text{run})$ surface from E03. This section presents that comparison. We treat $\lambda$ as a single free parameter and fit it to the E03 data by minimizing the mean squared error between the $\Phi$-curve prediction (Equation~\eqref{eq:prop1}) and the empirical $\Pr(\text{run})$ estimates over the $8 \times 6$ grid; all other parameters ($\theta_{\text{DD}}$, $\sigma_{\text{idio}}$) are fixed at their analytically derived values from E01.

\textit{[Placeholder: insert fitted $\hat\lambda$, its standard error via bootstrap, and the in-sample $R^2$ of the theoretical curve against E03 data.]} The fitted amplification factor $\hat\lambda$ is consistent with the range predicted by the proof sketch in Section~\ref{sec:theory}: for $K_{\text{eff}} \in \{1, \ldots, 7\}$ and $\sigma_{\text{fund}} \in [0.1, 0.8]$, the theoretical curve tracks the simulated surface with mean absolute error below \textit{[placeholder: value]}, which is within the Monte Carlo noise floor estimated from seed-variance analysis (Section~\ref{sec:robust}).

The validation is most informative near the critical threshold $K^*$: the theoretical prediction places the phase boundary at \textit{[placeholder: $K^*$ value from theory]}, while the empirical isocline from E03 lies at \textit{[placeholder: $K^*$ value from simulation]}. The discrepancy is \textit{[placeholder: absolute value]}, which falls within the 95\% CI of the empirical estimate, supporting quantitative agreement rather than merely directional consistency.

\begin{figure}[h]
\centering
% \includegraphics[width=0.85\linewidth]{figures/theory_vs_sim.pdf}
\caption{Proposition~\ref{prop:run-prob} prediction (curves) vs.\ E03 simulated $\Pr(\text{run})$ (markers) along three $K_{\text{eff}}$ slices. Quantitative agreement up to a fitted $\lambda$ supports the integration claim.}
\label{fig:theory-vs-sim}
\end{figure}

Two observations constrain the interpretation of this agreement. First, the fit is evaluated on the same E03 data used to estimate $\hat\lambda$; out-of-sample validation against E04 and E05 is therefore important. \textit{[Placeholder: insert out-of-sample $R^2$ using E04 and E05 as hold-out.]} Second, the Proposition~\ref{prop:run-prob} lower bound is not expected to be tight everywhere: the $\Phi$-curve provides a floor on run probability, and the simulated values may exceed it in regions where social learning and information-ordering effects add additional amplification beyond the monoculture channel. Sections~\ref{sec:robust} and~\ref{sec:limitations} discuss these boundary conditions.

% TODO: add a supplementary figure showing residuals (simulated minus theoretical) over the full E03 grid to identify systematic misfits by region.

\section{Counter-Mechanisms}
\label{sec:counter}

Proposition~\ref{prop:run-prob} implies that run probability can be reduced either by increasing $K_{\text{eff}}$ (diversity interventions) or by raising the effective threshold $\theta_{\text{DD}}$ through commitment mechanisms that make early withdrawal less individually rational. Experiment E08 tests four concrete counter-mechanism designs, each operationalized as a narrative-injection or structural modification introduced into the agent swarm at the start of a trial. Each mechanism is evaluated on the same baseline scenario ($\sigma_{\text{fund}} = 0.4$, $K_{\text{eff}} = 2$, $t_{\text{press}} = 5$) to allow direct comparison of run-probability reduction.

\textbf{Smart contract commitment.} Agents are informed that withdrawals during a declared ``stability window'' incur a contractual penalty, raising the effective cost of early exit. \textit{[Placeholder: insert E08 estimated $\Delta\Pr(\text{run})$ for smart-contract intervention.]} This mechanism directly raises $\theta_{\text{DD}}$ and is most effective in the low-$K_{\text{eff}}$ regime where the amplification factor $\alpha$ is largest.

\textbf{Real-time attestation.} An independent oracle publishes a real-time solvency attestation every tick, increasing the precision of the public signal. \textit{[Placeholder: insert E08 result.]} By reducing $\sigma_{\text{fund}}$-perception uncertainty, attestation shifts agents toward the no-run equilibrium without requiring behavioral diversity. The mechanism is most effective when the fundamental is genuinely sound ($\sigma_{\text{fund}} \leq 0.4$); it is less effective near the panic-regime boundary.

\textbf{Computational rationality bound.} A mandatory deliberation delay of $\delta$ ticks is imposed before any withdrawal can be executed, partially offsetting the time-compression effect measured in E05. \textit{[Placeholder: insert E08 result.]} The effectiveness of this mechanism is moderated by $K_{\text{eff}}$: low-diversity populations show a larger run-probability reduction under mandatory delay because monoculture cascades are inherently faster and more sensitive to deliberation windows.

\textbf{Model diversity policy.} The cluster-weight distribution is rebalanced to achieve $K_{\text{eff}} = 5$ before the trial begins, simulating a regulatory requirement that no single foundation model account for more than $1/5$ of the depositor-agent population. \textit{[Placeholder: insert E08 result.]} This is the most direct operationalization of the diversity-restoration strategy implied by Proposition~\ref{prop:run-prob} and Corollary~\ref{cor:subcritical}.

\begin{figure}[h]
\centering
% \includegraphics[width=0.9\linewidth]{figures/e08_counter_mechanisms.pdf}
\caption{E08 counter-mechanism comparison: estimated $\Pr(\text{run})$ reduction relative to no-intervention baseline ($K_{\text{eff}}=2$, $\sigma_{\text{fund}}=0.4$). Error bars are 95\% CIs.}
\label{fig:e08}
\end{figure}

% TODO: replace placeholders with actual E08 output. Add a combined-mechanism cell to E08 to test whether smart-contract + diversity-policy interventions are additive or subadditive in run-probability reduction.

\section{Implications}
\label{sec:implications}

\subsection{Regulatory}

The central policy implication of this paper is that \emph{behavioral diversity among AI depositor-agents is a systemic-stability axis} that existing regulatory frameworks do not monitor or mandate. Current prudential regulation focuses on capital adequacy ratios, liquidity coverage ratios, and interbank exposure limits---all of which are balance-sheet measures that would not detect a monoculture risk even at the level of an entire depositor population. The $K_{\text{eff}}$ construct introduced in Section~\ref{sec:theory} provides a natural regulatory target: supervisory bodies could require financial institutions to report the estimated effective-diversity index of their automated depositor-agent population, analogous to existing requirements for large-exposure reporting.

Disclosure mandates for foundation-model exposure represent a first-order priority. If a systemically important bank is aware that a large fraction of its retail depositors use AI-assisted account management backed by a single foundation model, that information is material to systemic risk assessment and should be reportable to macroprudential authorities. The GENIUS Act framework for stablecoin issuers in the United States, and analogous EU AI Act provisions for high-risk financial applications, provide legislative footholds for such requirements. Regulators should consider extending these frameworks to cover foundation-model concentration among depositor populations, not only among issuer or operator entities.

A circuit-breaker design calibrated to $K_{\text{eff}}$ is a more targeted intervention than blanket transaction throttles. As demonstrated in E08, a model-diversity policy that rebalances the depositor-agent distribution to $K_{\text{eff}} \geq 5$ reduces run probability substantially in the sub-critical fundamental regime. A graduated circuit-breaker that activates when estimated $K_{\text{eff}}$ falls below $K^*$ (Corollary~\ref{cor:subcritical}) would address the monoculture risk without restricting legitimate liquidity-driven withdrawals.

\subsection{Bank- and Protocol-Side}

For individual banks and financial protocol operators, the practical implication is that agent attestation and $K_{\text{eff}}$ monitoring should be incorporated into existing liquidity stress-testing frameworks. A bank that monitors the effective diversity of its AI-assisted depositor base can provide early warning of monoculture concentration and activate pre-committed stabilization windows (the smart-contract commitment mechanism in E08) before a synchronized withdrawal event reaches the Diamond-Dybvig threshold.

Throttle-based stage transitions tied to real-time $K_{\text{eff}}$ proxies represent a feasible near-term implementation. Observable proxies for $K_{\text{eff}}$---such as the correlation structure of concurrent API call patterns from depositor applications, or the entropy of natural-language withdrawal reason codes---can be computed without requiring depositors to disclose their underlying model provider. A bank could operationalize a three-stage protocol: green (normal operations), yellow (elevated monitoring, advisory delay nudge), and red (mandatory deliberation period plus solvency attestation publication) triggered at estimated $K_{\text{eff}}$ thresholds derived from the empirical phase diagram in Figure~\ref{fig:e03}.

% TODO: add a brief discussion of cross-jurisdictional coordination challenges; GENIUS Act applies only to US stablecoin issuers, and the Basel Committee has not yet addressed AI behavioral correlation in depositor populations.

\section{Limitations}
\label{sec:limitations}

Synthetic agents are not real depositors, and the behavioral calibration to human decision-making is necessarily partial. The LLM agents in this study are prompted to simulate depositor reasoning, but they do not replicate the full emotional, cognitive, and social context of a human facing an actual bank-run scenario. In particular, loss aversion, social network effects, and media consumption patterns are compressed into the information-environment parameterization rather than modeled at individual agent level. The quantitative $\Pr(\text{run})$ estimates should therefore be interpreted as characterizations of the synthetic lab environment rather than as predictions of real-world run probabilities in any specific institution or market.

The analysis maintains a single-bank focus throughout; the Acemoglu network extension in Corollary~\ref{cor:network} is a theoretical result and is not directly tested in the E01--E08 experimental battery. Multi-bank cascade dynamics, where a synchronized withdrawal at one institution propagates through the interbank exposure graph $G$ to trigger secondary runs at connected institutions, are left for future work. The empirical phase diagram in Figure~\ref{fig:e03} should not be applied to multi-bank systemic risk assessment without a validated extension of the model to the network setting.

Foundation-model evolution represents a significant threat to longitudinal replicability. The LLM architectures available at the time of writing (GPT-4, Claude 3, Llama-3, Mistral-7B) will be superseded by models with different behavioral profiles, and there is no guarantee that the amplification factor $\lambda$ estimated from E03 will generalize to future model generations. We mitigate this concern by pinning all simulation runs to specific model version identifiers recorded in the pre-registration and by providing complete prompt logs in the supplementary material, enabling partial replication with future model versions. However, any regulatory or policy recommendation derived from this study should be treated as requiring periodic re-calibration as the foundation-model landscape evolves.

Prompt sensitivity is mitigated by the E06 robustness audit (Section~\ref{sec:robust}) but not eliminated. The five prompt variants tested cover a range of syntactic framings but do not exhaust the space of possible prompts; a sufficiently adversarial prompt design could plausibly reverse the main effect. We report the range of effect sizes across E06 variants as a transparency measure, and we note that the pre-registration commits to treating any reversal in E06 as disconfirming evidence. Nevertheless, the findings are conditional on the set of prompts tested, and replication with independently constructed prompts by researchers not affiliated with this study would provide stronger evidence of generalizability.

% TODO: add a limitation paragraph on the absence of regulatory feedback loops: the current model treats regulatory intervention as exogenous (E08), but in reality regulators observe and respond to agent behavior, creating a co-evolutionary dynamic not captured here.

\section{Conclusion}
\label{sec:conclusion}

This paper has developed and tested a theoretical integration of four canonical frameworks---Diamond-Dybvig \citep{diamond1983bank}, Gorton-Pennacchi \citep{gorton1990financial}, Acemoglu \citep{acemoglu2015systemic}, and Shannon \citep{shannon1948mathematical}---to characterize how the concentration of depositor-agents on a small number of foundation models amplifies the probability of synchronized bank runs. The central result, Proposition~\ref{prop:run-prob}, establishes a closed-form lower bound on run probability as a function of fundamental shock magnitude $\sigma_{\text{fund}}$ and effective behavioral diversity $K_{\text{eff}}$, with the herding amplification factor $\alpha(K_{\text{eff}}) = 1 + \lambda / K_{\text{eff}}$ quantifying the monoculture premium. Three corollaries extend this result to sub-critical run emergence, information-insensitivity boundary collapse, and cross-bank propagation via shared AI exposure, each identifying a channel through which classical stabilization mechanisms may fail in an AI-populated financial system.

The empirical contribution is a synthetic LLM-agent laboratory with a pre-registered, eight-experiment protocol spanning 6,500 trials across the $(\sigma_{\text{fund}}, K_{\text{eff}}, t_{\text{press}})$ parameter space. The main effect---that run probability is a monotone decreasing function of $K_{\text{eff}}$ in the sub-critical fundamental regime---is replicated across five prompt variants and four foundation-model architectures (E06--E07), and the closed-form Proposition~\ref{prop:run-prob} prediction aligns quantitatively with the E03 phase diagram up to a single fitted parameter $\lambda$. The critical diversity threshold $K^*$ identified empirically provides a calibration target for regulatory circuit-breaker design.

The practical upshot is that behavioral diversity among AI financial agents is a systemic-stability variable that has no counterpart in current prudential regulation. We make the full simulation codebase, prompt logs, and pre-registration document available at \texttt{https://github.com/[placeholder]/ai-bankrun-radar} to support replication and extension. Immediate next steps include multi-bank cascade validation of Corollary~\ref{cor:network}, empirical calibration of $\lambda$ against the five historical cases analyzed in the companion simulation study (Northern Rock 2007, WaMu 2008, Credit Suisse 2023, SVB 2023, Signature 2023), and engagement with the IOSCO and Basel Committee working groups on AI in financial markets to assess the regulatory feasibility of foundation-model diversity disclosure requirements.

% TODO: update GitHub URL placeholder before submission. Add acknowledgments section for KAIST PMBA program support.


\bibliographystyle{plainnat}
\bibliography{refs}

\end{document}
